% GENERATED FILE. Edit canonical lessons, then run npm run generate:books.
% canonical-source-hash: fnv1a64:d25ffb67828e1092
% canonical-lessons: ZH-C464-fahuo, ZH-C464-huoqi, ZH-C464-qiren, ZH-C464-chuqi, ZH-C464-fennu

\chapter{Losing My Temper}
\label{ch:zh-losing-my-temper}
\hlchaptermodality{464}

\begin{chapteropening}
\textbf{By the end of this chapter:} \emph{I can say someone lost their temper or has a hot temper, that someone is infuriating, and that I am furious.}

Complete the last lesson of chapter 464: I can say someone lost their temper or has a hot temper, that someone is infuriating, and that I am furious.
\end{chapteropening}

\section[\texorpdfstring{\zh{发火} (fāhuǒ)}{发火 (fāhuǒ)}]{\zh{发火} (fāhuǒ) — to lose one's temper}

\label{lesson:ZH-C464-fahuo}



\begin{quote}
Before the new one: say the Mandarin for to tremble, then the Mandarin for a nightmare.
\end{quote}

\subsection*{You'll want to know: \zh{发火}}
\textbf{\zh{发火}} — \emph{fāhuǒ} — \textquotedblleft{}to lose one's temper\textquotedblright{}.

\begin{quote}
\textbf{\zh{他发火的时候我很怕他}} \emph{tā fāhuǒ de shíhou wǒ hěn pà tā} — I'm really scared of him when he loses his temper.
\end{quote}

\begin{cousinweb}
\textbf{\zh{发火}} (\emph{fāhuǒ}) is built from \zh{发} 'to send out' and \zh{火} 'fire', both inherited from Middle Chinese, \emph{pjot} and \emph{xwaX} (Guangyun, in Baxter's transcription), which glosses the first 'to start' and 'to raise'. Anger as a fire that breaks out is the usual explanation, not one the dictionary gives.
\end{cousinweb}

\begin{grammarlens}[title={What you've built}]
One more word for saying how it was.
\end{grammarlens}

\subsection*{Your turn}
\begin{itemize}
\raggedright
  \item \emph{Say it:} \emph{fāhuǒ}
  \item \emph{Say it:} \emph{fāhuǒ}, once more
  \item \emph{Say it:} \emph{tā fāhuǒ de shíhou wǒ hěn pà tā} — I'm really scared of him when he loses his temper.
  \item \emph{Recall:} say the Mandarin for to tremble, then the Mandarin for a nightmare, then say \emph{fāhuǒ} again
\end{itemize}

\begin{tcolorbox}[breakable,colback=teal!4,colframe=teal!35!black,title={Before you move on}]
What does \zh{发火} mean? (\textquotedblleft{}To lose one's temper\textquotedblright{}.) Say it once more.
\end{tcolorbox}
\section[\texorpdfstring{\zh{火气} (huǒqì)}{火气 (huǒqì)}]{\zh{火气} (huǒqì) — temper; anger}

\label{lesson:ZH-C464-huoqi}



\begin{quote}
Before the new one: say the Mandarin for a nightmare, then the Mandarin for to lose one's temper.
\end{quote}

\subsection*{You'll want to know: \zh{火气}}
\textbf{\zh{火气}} — \emph{huǒqì} — \textquotedblleft{}temper; anger\textquotedblright{}.

\begin{quote}
\textbf{\zh{他的火气比我大}} \emph{tā de huǒqì bǐ wǒ dà} — He has a worse temper than I do.
\end{quote}

\begin{cousinweb}
\textbf{\zh{火气}} (\emph{huǒqì}) is built from \zh{火} 'fire' and \zh{气} 'breath', both inherited from Middle Chinese, \emph{xwaX} and \emph{khj+jH} (Guangyun, in Baxter's transcription), which glosses the second 'breath'. The sense 'anger' is not given there.
\end{cousinweb}

\begin{grammarlens}[title={What you've built}]
One more word for saying how it was.
\end{grammarlens}

\subsection*{Your turn}
\begin{itemize}
\raggedright
  \item \emph{Say it:} \emph{huǒqì}
  \item \emph{Say it:} \emph{huǒqì}, once more
  \item \emph{Say it:} \emph{tā de huǒqì bǐ wǒ dà} — He has a worse temper than I do.
  \item \emph{Recall:} say the Mandarin for a nightmare, then the Mandarin for to lose one's temper, then say \emph{huǒqì} again
\end{itemize}

\begin{tcolorbox}[breakable,colback=teal!4,colframe=teal!35!black,title={Before you move on}]
What does \zh{火气} mean? (\textquotedblleft{}Temper; anger\textquotedblright{}.) Say it once more.
\end{tcolorbox}
\section[\texorpdfstring{\zh{气人} (qìrén)}{气人 (qìrén)}]{\zh{气人} (qìrén) — to be infuriating; to make people angry}

\label{lesson:ZH-C464-qiren}



\begin{quote}
Before the new one: say the Mandarin for to lose one's temper, then the Mandarin for a hot temper.
\end{quote}

\subsection*{You'll want to know: \zh{气人}}
\textbf{\zh{气人}} — \emph{qìrén} — \textquotedblleft{}to be infuriating; to make people angry\textquotedblright{}.

\begin{quote}
\textbf{\zh{这个人真气人}} \emph{zhège rén zhēn qìrén} — That man is really infuriating.
\end{quote}

\begin{cousinweb}
\textbf{\zh{气人}} (\emph{qìrén}) is built from \zh{气} 'to anger' and \zh{人} 'a person', both inherited from Middle Chinese, \emph{khj+jH} and \emph{nyin} (Guangyun, in Baxter's transcription), which glosses the first 'breath'; the sense 'to anger' is not given there.
\end{cousinweb}

\begin{grammarlens}[title={What you've built}]
One more word for saying how it was.
\end{grammarlens}

\subsection*{Your turn}
\begin{itemize}
\raggedright
  \item \emph{Say it:} \emph{qìrén}
  \item \emph{Say it:} \emph{qìrén}, once more
  \item \emph{Say it:} \emph{zhège rén zhēn qìrén} — That man is really infuriating.
  \item \emph{Recall:} say the Mandarin for to lose one's temper, then the Mandarin for a hot temper, then say \emph{qìrén} again
\end{itemize}

\begin{tcolorbox}[breakable,colback=teal!4,colframe=teal!35!black,title={Before you move on}]
What does \zh{气人} mean? (\textquotedblleft{}To be infuriating; to make people angry\textquotedblright{}.) Say it once more.
\end{tcolorbox}
\section[\texorpdfstring{\zh{出气} (chūqì)}{出气 (chūqì)}]{\zh{出气} (chūqì) — to vent one's anger; to take it out on someone}

\label{lesson:ZH-C464-chuqi}



\begin{quote}
Before the new one: say the Mandarin for a hot temper, then the Mandarin for to be infuriating.
\end{quote}

\subsection*{You'll want to know: \zh{出气}}
\textbf{\zh{出气}} — \emph{chūqì} — \textquotedblleft{}to vent one's anger; to take it out on someone\textquotedblright{}.

\begin{quote}
\emph{nǐ shēngqì de shíhou bié ná wǒ chūqì} — When you get angry, don't take it out on me.
\end{quote}

\begin{cousinweb}
\textbf{\zh{出气}} (\emph{chūqì}) is built from \zh{出} 'out' and \zh{气} 'breath', both inherited from Middle Chinese, \emph{tsyhwit} and \emph{khj+jH} (Guangyun, in Baxter's transcription), which glosses the second 'breath'; letting angry breath out is the usual explanation, not one the dictionary gives.
\end{cousinweb}

\begin{grammarlens}[title={What you've built}]
One more word for saying how it was.
\end{grammarlens}

\subsection*{Your turn}
\begin{itemize}
\raggedright
  \item \emph{Say it:} \emph{chūqì}
  \item \emph{Say it:} \emph{chūqì}, once more
  \item \emph{Say it:} \emph{nǐ shēngqì de shíhou bié ná wǒ chūqì} — When you get angry, don't take it out on me.
  \item \emph{Recall:} say the Mandarin for a hot temper, then the Mandarin for to be infuriating, then say \emph{chūqì} again
\end{itemize}

\begin{tcolorbox}[breakable,colback=teal!4,colframe=teal!35!black,title={Before you move on}]
What does \zh{出气} mean? (\textquotedblleft{}To vent one's anger; to take it out on someone\textquotedblright{}.) Say it once more.
\end{tcolorbox}
\section[fènnù]{fènnù — furious; indignant}

\label{lesson:ZH-C464-fennu}



\begin{quote}
Before the new one: say the Mandarin for to be infuriating, then the Mandarin for to take it out on someone.
\end{quote}

\subsection*{You'll want to know: fènnù}
\emph{fènnù} — \textquotedblleft{}furious; indignant\textquotedblright{}.

\begin{quote}
\emph{tīng le tā de huà wǒ gǎnjué hěn fènnù} — When I heard what he said, I was furious.
\end{quote}

Its characters come later; for now it is a word for the ear.

\begin{cousinweb}
\emph{fènnù} is built from \emph{fèn} and \emph{nù}, both inherited from Middle Chinese, \emph{bjunX} and \emph{nuH} (Guangyun, in Baxter's transcription), which glosses the first 'vexed' and the second 'anger'.
\end{cousinweb}

\begin{grammarlens}[title={What you've built}]
One more word for saying how it was.
\end{grammarlens}

\subsection*{Your turn}
\begin{itemize}
\raggedright
  \item \emph{Say it:} \emph{fènnù}
  \item \emph{Say it:} \emph{fènnù}, once more
  \item \emph{Say it:} \emph{tīng le tā de huà wǒ gǎnjué hěn fènnù} — When I heard what he said, I was furious.
  \item \emph{Recall:} say the Mandarin for to be infuriating, then the Mandarin for to take it out on someone, then say \emph{fènnù} again
\end{itemize}

\begin{tcolorbox}[breakable,colback=teal!4,colframe=teal!35!black,title={Before you move on}]
What does fènnù mean? (\textquotedblleft{}Furious; indignant\textquotedblright{}.) Say it once more.
\end{tcolorbox}
